\section{Features force companions}\label{sec:companions}

% ledger: A1-A9, A17, O6; source proof/ALIGNMENT_PROOF.md A§1-A§4 (A supersedes R where both treat a step)
This section and the next prove that every tiling by $\solid$ is registered
(\cref{fig:companion}). The argument is
continuous geometry: the pyramids are tiny, but their edges carry dihedral angles that only a
pyramid of the same magnitude and opposite polarity can complete, and that forces every feature
of every tile to be matched, whole, by one feature of one neighbour. We follow the written
alignment proof of the repository, whose lemma labels (A-L2.1 and so on) we give with each
statement together with the Lean theorem that proves it and that theorem's tier
(\cref{tab:hypotheses}); the written proofs are the exposition, and where a formal proof argues
differently we say so after the written one.

% notation; source A§1
Fix a feature $F$ of a tile and write $E(F)$ for its \emph{closed edge graph}: the four edges of
the base square and the four ridges from the base corners to the apex, a compact connected graph
without isolated vertices. In tangent coordinates $u$ on the panel and outward normal coordinate
$z$, the feature of signed height $h=a/10000$ is the graph of the tent function
$f_h(u)=h\max(0,1-\|u\|_\infty/\eta)$, and near it the solid is the hypograph $z\le f_h(u)$.
Throughout, a \emph{packing} is a family of isometric copies of $\solid$ with pairwise disjoint
interiors, and a tiling is a packing whose union is $\R^3$.

\begin{lemma}[local finiteness; A-L2.1, \tier{T1} \lean{local\_finiteness}]\label[lemma]{lem:locfin}
Every packing by congruent copies of $\solid$ is locally finite.
\end{lemma}

\begin{proof}
Choose a ball of positive radius inside the interior of $\solid$ and carry it along with each
tile. The carried balls have pairwise disjoint interiors. A tile meeting a bounded set $K$ has its
ball inside a fixed bounded enlargement of $K$, since $\solid$ has bounded diameter, and only
finitely many disjoint balls of equal positive volume fit there.
\end{proof}

% ledger: A2; Lean cone_sector_budgets_holds (T1n)
\begin{lemma}[cone budgets; A-L2.2, written; Lean \tier{T1n} \lean{cone\_sector\_budgets\_holds}]\label[lemma]{lem:cones}
At a point common to finitely many tiles of a packing, the interiors of their tangent cones are
disjoint, so their solid angles sum to at most $4\pi$, with equality in a tiling. At a generic
point of an edge, a perpendicular cross-section is partitioned into planar sectors whose angles
sum to $2\pi$ in a tiling.
\end{lemma}

\begin{proof}
Near a boundary point a polyhedral set agrees with a translate of its tangent cone, and one
radius serves the finitely many incident tiles; a common open direction would give overlapping
interiors. Cone boundaries are finite unions of planar pieces and have spherical area zero.
Coverage gives equality in a tiling. The planar statement is the same argument in the
perpendicular plane.
\end{proof}

The word \emph{generic} is made precise as follows. Fix a closed root edge and the finitely many
tiles meeting a small neighbourhood of it. Remove their vertices on the edge, their intersections
with non-collinear edges, and the intersections with face planes that meet the root line in a
single point: finitely many points. At every remaining point an incident tile either has a face
containing the line, which contributes a half-plane sector of angle $\pi$, or an edge collinear
with it, which contributes an ordinary dihedral sector. No tile can contain such a point in its
interior, since $\solid$ is the closure of its interior and the root has interior points nearby.
Nothing here assumes that a neighbouring facet coincides with a facet of the root; contacts that
are not face-to-face are handled by the same sectors.

% ledger: A3, O6, A17 (E3); Lean deviations_distinct (T1), mesh_angle_audit (T1n)
Put $t_j=j/100$ for $j=1,\dots,12$, and
\[
  \alpha_j=\arctan t_j,\qquad \beta_j=\arccos\frac{1}{1+t_j^2}.
\]

\begin{lemma}[complete dihedral list; A-L3.1, written; Lean \tier{T1n}
\lean{complete\_dihedral\_list\_holds}; its arithmetic and the mesh audit are separate Lean
theorems, see below]\label[lemma]{lem:dihedral}
A protrusion of magnitude $j$ has base dihedral $\pi+\alpha_j$ and ridge dihedral $\pi-\beta_j$;
a recess has the complementary angles. Every other edge of $\solid$ has dihedral $\pi/2$ or
$3\pi/2$, and flat triangulation edges have angle $\pi$. All $24$ deviations $\alpha_j,\beta_j$
are distinct and less than $\pi/4$.
\end{lemma}

\begin{proof}
The facets of the tent have slopes $(0,0)$, $(\pm t_j,0)$ and $(0,\pm t_j)$; their outward normals
give the formulas and the signs. The tube construction changes no carrier edge, and the disjoint
feature supports create no other kind of edge. Both sequences increase strictly, and a
coincidence $\alpha_i=\beta_j$ would force $1+(i/100)^2=(1+(j/100)^2)^2$, that is
$10000\,i^2=20000\,j^2+j^4$; then $10\mid j$, so $j=10$ and $i^2=201$, impossible. The bound
$\pi/4$ follows from $t_j<1$ and $(1+t_j^2)^2<2$ at $t_{12}=3/25$.
\end{proof}

The unsolvability of $10000\,i^2=20000\,j^2+j^4$ over $1\le i,j\le12$ is a Lean theorem by kernel
\lean{decide} (\lean{deviations\_distinct}). The lemma is about the ideal construction; a separate
audit, a Lean theorem using \lean{native\_decide} (\lean{mesh\_angle\_audit}), classifies every
one of the $6{,}408$ edges of the actual mesh by exact rational outward normals, finds exactly
$1{,}536$ feature edges with $32$ of each sign for each of the $24$ deviations, and checks that
every other edge is flat or has orthogonal normals (\cref{tab:mesh}).

% ledger: A4, A17 (E3); Lean generic_feature_partner_holds (T1n)
\begin{lemma}[one complementary partner; A-L3.2, written; Lean \tier{T1n}
\lean{generic\_feature\_partner\_holds}]\label[lemma]{lem:partner}
At a generic point of a feature edge in any tiling exactly two tiles are incident: the root, and
one tile whose incident edge has the same type (base or ridge), the same magnitude, and the
complementary dihedral angle.
\end{lemma}

\begin{proof}
Every feature sector lies strictly between $3\pi/4$ and $5\pi/4$. Three feature sectors exceed
$2\pi$, as do two feature sectors together with any ordinary sector, whose angle is at least
$\pi/2$. A single feature sector is impossible: its angle is $\pi$ plus a nonzero deviation of
magnitude less than $\pi/4$, while the remaining sectors are multiples of $\pi/2$ (a face
containing the line contributes $\pi$, as noted above), and such a sum is not $2\pi$. The root
supplies one feature sector, so there are exactly two and nothing else; their deviations are
opposite and equal in magnitude, and \cref{lem:dihedral} identifies the type and the magnitude.
\end{proof}

This does not yet say that an edge angle forces a whole feature. That needs a bound at the
vertices of the feature, where sectors are not available.

% ledger: A5, A6; Lean feature_circular_cone_containment_holds (T1n), circular_cone_solid_angle_holds (T1)
\begin{lemma}[solid-angle bound on the feature graph; A-L4.1, written; Lean \tier{T1n}
\lean{feature\_circular\_cone\_containment\_holds}, with the cone formula and
its arithmetic \tier{T1} \lean{circular\_cone\_solid\_angle\_holds}]\label[lemma]{lem:solidangle}
At every point of $E(F)$, including the apex and the base corners, the interior solid angle of
the tile exceeds $4\pi/3$; for $\solid$ it exceeds $7\pi/4$.
\end{lemma}

\begin{proof}
The tent is Lipschitz with constant at most $L=3/25$, on its faces, along its creases, and across
the transition to the flat panel. At any point of the graph the tile's interior tangent cone
contains the circular cone $\{(v,w): w<-L\|v\|_2\}$, whose solid angle is
$\Omega(L)=2\pi\bigl(1-L/\sqrt{1+L^2}\bigr)$. This exceeds $4\pi/3$ exactly when $8L^2<1$, and here
$8L^2=72/625$; it exceeds $7\pi/4$ when $63L^2<1$, and $63L^2=567/625$.
\end{proof}

Only the weaker bound is used below: by \cref{lem:cones}, three distinct tiles cannot share a
point lying on a feature graph of each of them.

\begin{figure}[ht]
  \centering
  \resizebox{\linewidth}{!}{% fig_companion_sectors.tex — Fig. 4.1: the sector budget at a generic feature-edge point (left),
% why three feature sectors cannot meet (middle), and the solid-angle cone at a vertex of the
% feature graph (right). Schematic. Used by sec4_companions.tex.
\begin{tikzpicture}[x=1cm,y=1cm,font=\small,>={Latex[length=1.6mm]}]
  \begin{scope}
    \fill[red!15] (0,0) -- (-2.2,0) arc (180:360:2.2) -- (2.2,0) -- (2.2*0.94,2.2*0.34) arc (20:0:2.2) -- cycle;
    \fill[blue!15] (0,0) -- (2.2*0.94,2.2*0.34) arc (20:180:2.2) -- cycle;
    \draw[thick] (-2.2,0) -- (2.2,0);
    \draw[thick] (0,0) -- (2.2*0.94,2.2*0.34);
    \node at (0,-1.1) {root: $\pi+\alpha_j$};
    \node at (-0.6,1.0) {partner: $\pi-\alpha_j$};
    \node[anchor=north,align=center,text width=5.2cm] at (0,-2.5)
      {a generic point of a feature edge: exactly two sectors, the same magnitude $j$, the same type, complementary angles};
  \end{scope}
  \begin{scope}[shift={(7.4,0)}]
    \draw[thick] (0,0) circle (1.6);
    \foreach \a in {0,120,240} { \draw[thick] (0,0) -- (\a:1.6); }
    \node at (60:1.0) {$>\tfrac{3\pi}{4}$};
    \node at (180:1.0) {$>\tfrac{3\pi}{4}$};
    \node at (300:1.0) {$>\tfrac{3\pi}{4}$};
    \node[anchor=north,align=center,text width=5.2cm] at (0,-2.5)
      {three feature sectors would exceed $2\pi$; two plus an ordinary sector ($\ge\pi/2$) too};
  \end{scope}
  \begin{scope}[shift={(14.8,0)}]
    \draw[thick] (-1.9,0) -- (1.9,0);
    \fill[black!8] (0,0) -- (-1.8,-0.22) -- (-1.8,-1.5) -- (1.8,-1.5) -- (1.8,-0.22) -- cycle;
    \draw[thick] (0,0) -- (-1.8,-0.22); \draw[thick] (0,0) -- (1.8,-0.22);
    \fill (0,0) circle (1.2pt);
    \node[anchor=south] at (0,0.08) {apex or base corner};
    \node[anchor=north,align=center,text width=5.2cm] at (0,-2.5)
      {interior solid angle $>4\pi/3$ at every point of $E(F)$: no third tile on the graph};
  \end{scope}
\end{tikzpicture}
}
  \caption{The companion argument: sectors at a generic edge point (\cref{lem:partner}) and the
  solid-angle cone at a vertex of the feature graph (\cref{lem:solidangle}).}
  \label{fig:companion}
\end{figure}

% ledger: A7; Lean connected_feature_companion_holds (T1n)
\begin{theorem}[one tile accompanies the whole feature graph; A-T4.2, written; Lean \tier{T1n}
\lean{connected\_feature\_companion\_holds}]\label{thm:companion}
For each feature $F$ of a tile $R$ in any tiling there is a unique other tile $S$ such that
$E(F)$ is contained in the union of the feature edge graphs of $S$; moreover $E(F)$ lies in one
feature edge graph of $S$.
\end{theorem}

\begin{proof}
By \cref{lem:locfin}, only finitely many tiles other than $R$ meet $E(F)$. For each such tile
$S\ne R$, let $C_S$ be the intersection of $E(F)$ with the union of the feature edge graphs of
$S$; each $C_S$ is closed in $E(F)$. The sets $C_S$ cover $E(F)$: the generic interior points of the root's edges are covered
by \cref{lem:partner}, they are dense in $E(F)$, and a finite union of closed sets containing a
dense subset is everything. The sets $C_S$ are pairwise disjoint: a point in $C_S\cap C_T$ with
$S\ne T$ would lie on a feature graph of $R$, of $S$ and of $T$, which \cref{lem:cones,lem:solidangle}
forbid. A connected space is not the union of two or more nonempty members of a finite disjoint
closed cover, each member being also open. Hence exactly one $C_S$ is nonempty and equals $E(F)$.
The feature graphs of $S$ are disjoint compact sets, and the same connectedness argument applied
to $E(F)$ inside their union puts $E(F)$ in one of them.
\end{proof}

This single argument disposes of every case in which a partner might change: at an interior edge
point, at a subdivision vertex of the mesh, at the apex, at a base corner, or between two
features of the same neighbour. It needs no face-to-face hypothesis and no continuation rule
supplied by the observer.

% ledger: A8, A17 (E1); Lean feature_containment_rigidity_holds (T1n; formal proof by compactness)
\begin{lemma}[containment implies equality; A-L4.3, written, with erratum E1; Lean \tier{T1n}
\lean{feature\_containment\_rigidity\_holds}]\label[lemma]{lem:equal}
The features $F$ and $F'$ of \cref{thm:companion} have the same closed edge graph, the same
magnitude, and opposite polarity; their triangulated surfaces coincide.
\end{lemma}

\begin{proof}
Choose a point in the interior of a base edge of $F$ that avoids the vertices of the companion's
graph and the finite exceptional set of \cref{lem:cones}; such a point exists. There
\cref{lem:partner} gives $F'$ the same magnitude $j$ and the same edge type. A feature of
magnitude $j$ has total edge length $\ell_j=8\eta+4\sqrt{2\eta^2+(j/10000)^2}$. We know $E(F)$ is
a closed subset of $E(F')$; if the inclusion were proper, a point of $E(F')$ outside $E(F)$ would
have positive distance from the closed set $E(F)$, and since $E(F')$ has no isolated vertices it
would contribute a nonzero subsegment outside $E(F)$, giving $\ell_j<\ell_j$ by additivity of
length. So the graphs are equal. In the graph the apex is the unique vertex of degree four and
the base corners have degree three, so equality recovers the base square, the apex, and the four
triangles; complementary base dihedrals force opposite polarity.
\end{proof}

The formal proof argues differently: the graph $E(F)$ is compact and the placement carrying it
into $E(F')$ is an isometry, and an isometric self-map of a compact metric space is onto, so the
containment is an equality without the length count; the premise that the two tiles are distinct
is used there.

% ledger: A9; Lean companion_pose_discrete_holds (T1n)
\begin{corollary}[a feature match discretizes the pose; A-C4.4, written; Lean \tier{T1n}
\lean{companion\_pose\_discrete\_holds}]\label[corollary]{cor:discrete}
Normalize the root's pose to $(I,0)$. Every feature companion has a signed permutation matrix $G$
as linear part and translation $t=c_u-Gc_v\in\tfrac18\Z^3$, where $c_u,c_v$ are native feature
centres with $a_u=-a_v$ and outward base normals $n_u=-Gn_v$.
\end{corollary}

\begin{proof}
A rigid motion matching two complete pyramid graphs carries base-edge directions and base-plane
normal to each other, and those are the coordinate directions, so its orthogonal part is a
signed permutation, of either determinant. The graph identifies the centre; opposite polarity and
a common apex force opposite outward normals; all centres lie in $\tfrac18\Z^3$.
\end{proof}

No global alignment has been concluded yet: only the poses of complete feature companions have
been reduced to a finite list, which the next section examines exhaustively. The mechanism is the
one Socolar and Taylor use for their three-dimensional tile, where plugs enforce a matching rule
by shape \cite{ST11}; here the pyramids enforce, by dihedral angles and solid angles alone, that
every feature is matched by a feature of the opposite polarity on one neighbour.
